\section{The Koenigs identification}
\label{a:sec:koenigs}

A dynamical argument explains why no elementary closed form should be expected. In
short: $A(p)$ is a value of the Koenigs linearising function of the logistic map,
and that function is known to be non-elementary throughout the parameter range
required.

\subsection{From the survival recursion to the logistic map}
\label{a:sec:tologistic}

The substitution has already been made in \cref{a:sec:product}. With $S_n = 2w_n$
and $r = 2p$, the survival recursion \cref{a:eq:SnIter} becomes the logistic
iteration
\begin{equation}
\label{a:eq:logisticagain}
w_{n+1} = f_r(w_n),
\qquad
f_r(w) = r\,w(1-w),
\qquad
w_0 = \tfrac12,
\end{equation}
with $r\in[0,1)$ throughout the subcritical regime, and
\begin{equation}
\label{a:eq:Afromw}
A(p) = 2\lim_{n\to\infty}\frac{w_n}{r^n} .
\end{equation}

\subsection{Koenigs linearisation}
\label{a:sec:koenigsdef}

\begin{definition}[Koenigs function]
\label{a:def:koenigs}
Let $f_r(z) = rz(1-z)$ have a fixed point at $z=0$ with multiplier
$r = f_r'(0)$ satisfying $0<|r|<1$. Koenigs \cite{Koenigs1884} proved that there
is a unique function $\psi_r$, holomorphic in a neighbourhood of the origin,
satisfying the Schröder equation
\begin{equation}
\label{a:eq:schroder}
\psi_r\bigl(f_r(z)\bigr) = r\,\psi_r(z),
\qquad
\psi_r(0)=0,
\qquad
\psi_r'(0)=1 .
\end{equation}
Equivalently,
\begin{equation}
\label{a:eq:koenigslimit}
\psi_r(z) = \lim_{n\to\infty}r^{-n}f_r^{\circ n}(z)
\end{equation}
uniformly near the origin. The normalisation $\psi_r'(0)=1$ is equivalent to
$\psi_r(z)/z\to1$ as $z\to0$.
\end{definition}

\noindent
The Koenigs function is a change of coordinate that straightens the dynamics: in
the $\psi$ coordinate the nonlinear map $f_r$ becomes multiplication by $r$
(\cref{a:fig:koenigs}). Here $f_r$ is exactly \cref{a:eq:logisticagain}, with
$z=w_n$ and $f_r(z)=w_{n+1}$, and the multiplier is $r=2p\in(0,1)$ --- an attracting
fixed point, which is the classical Koenigs setting. Standard accounts are Milnor
\cite{Milnor2006} and Carleson and Gamelin \cite{Carleson1993}.

\begin{figure}[ht]
\centering
\includegraphics[width=0.72\textwidth]{figure2_koenigs_linearization.png}
\caption{Koenigs linearisation of the logistic map near an attracting fixed point.
The nonlinear dynamics $x\mapsto f_r(x)=rx(1-x)$ are conjugated by $\psi_r$ to the
linear map $y\mapsto ry$; the diagram displays the functional relation
$\psi_r(f_r(x)) = r\psi_r(x)$ and the role of $\psi_r$ as the coordinate change
that straightens the dynamics.}
\label{a:fig:koenigs}
\end{figure}

\subsection{The identity}
\label{a:sec:identity}

Koenigs' theorem is local: the series \cref{a:eq:koenigslimit} is guaranteed to converge only
on a disc about the origin, and for the logistic map that disc has radius
$(1-r)/r$. The evaluation point of interest is $w_0=\tfrac12$, and
$\tfrac12<(1-r)/r$ fails as soon as $r\ge\tfrac23$, that is $p\ge\tfrac13$. For
those parameters $\tfrac12$ lies outside the germ, and writing
``$\psi_r(\tfrac12)$'' as a naive evaluation of a power series would be
meaningless.

The remedy is the functional equation itself. For $r\in(0,1)$ the origin is
attracting on the real line and the whole interval $(0,1)$ lies in its real basin,
so the orbit $w_n = f_r^{\circ n}(w_0)$ enters the germ after finitely many steps.
Define the \emph{basin extension}
\begin{equation}
\label{a:eq:basinkoenigs}
\Psi_r(w) := \lim_{n\to\infty}r^{-n}f_r^{\circ n}(w),
\qquad w\in(0,1) ,
\end{equation}
equivalently $\Psi_r(w) = r^{-N}\psi_r\bigl(f_r^{\circ N}(w)\bigr)$ for any $N$
large enough that $f_r^{\circ N}(w)$ lies in the germ. The two definitions agree
where both apply, because \cref{a:eq:schroder} makes the right-hand side
independent of $N$, and $\Psi_r$ is the unique analytic continuation of $\psi_r$
along the orbit.

\begin{proposition}[The constant as a Koenigs value]
\label{a:prop:AisKoenigs}
For $p\in(0,\tfrac12)$ and $r=2p$,
\begin{equation}
\label{a:eq:AeqKoenigs}
A(p) = 2\,\Psi_r\!\lrb{\tfrac12}
= 2\lim_{n\to\infty}r^{-n}f_r^{\circ n}\!\lrb{\tfrac12} .
\end{equation}
\end{proposition}

\begin{proof}
Iterating \cref{a:eq:schroder} along the orbit $w_0,w_1,w_2,\dots$ of
\cref{a:eq:logisticagain}, from any point at which the germ applies, gives
$\psi_r(w_{n}) = r\,\psi_r(w_{n-1})$ and hence
\begin{equation}
\label{a:eq:induction}
\Psi_r(w_0) = \frac{\psi_r(w_n)}{r^n}
\qquad\text{for every $n$ large enough that $w_n$ lies in the germ.}
\end{equation}
Passing to the limit, $\Psi_r(w_0) = \lim_n\psi_r(w_n)/r^n$. The normalisation
$\psi_r(0)=0$, $\psi_r'(0)=1$ means $\psi_r(z)/z\to1$ as $z\to0$, and
$w_n = S_n/2\to0$ because $r<1$; replacing $\psi_r(w_n)$ by $w_n$ in the limit is
therefore legitimate, and $\Psi_r(w_0) = \lim_nw_n/r^n$. Comparing with
\cref{a:eq:Afromw} and using $w_0 = S_0/2 = \tfrac12$ gives
$A(p) = 2\Psi_r(\tfrac12)$.
\end{proof}

We write $A(p) = 2\psi_r(\tfrac12)$ below as shorthand for this basin value.

\begin{remark}[Germ versus basin]
\label{a:rem:basin}
Differential algebraicity is a local
condition preserved under analytic continuation: if an algebraic differential
equation held for $\psi_r$ on any subdomain, it would persist along the
continuation defining $\Psi_r$, and conversely. Hypertranscendence of the germ and
of the basin extension are therefore equivalent statements, and the identity
\cref{a:eq:AeqKoenigs} is unaffected by which one is under discussion.
\end{remark}

\begin{figure}[ht]
\centering
\includegraphics[width=0.66\textwidth]{figure3_numerical_koenigs.png}
\caption{Numerical approximation of a Koenigs function for a logistic multiplier
in $(0,1)$, with low-order Taylor polynomials near the origin. Elementary closed
forms exist only for exceptional multipliers outside this range; those cases are
derived in \cref{a:app:integrable}.}
\label{a:fig:koenigsnumeric}
\end{figure}
